% ============================================================================
% Fichier  : partie0/P00_quarante_cinq_questions.tex
% Titre    : Prologue -- Quarante-Cinq Questions et Aucune Réponse Fiable
% Auteur   : MINKA MI NGUIDJOÏ Thierry Emmanuel
% Version  : 1.0 -- octobre 2026
% ============================================================================
% Source du cas : illustration LizbiZ de l'article A0, « Vers l'auditologie »
% (HAL hal-05701007), section 4. Les faits repris de l'article (grille de 45
% éléments en 5 domaines, domaine D incohérent, H ~ 1,2 bit, décomposition
% 5 x 5, usage stratégique de NEK, régression de 35 % à 18 mois, résultats à
% six mois) ne doivent pas être modifiés. Les autres faits sont propres au
% cours et compatibles avec l'article.
% ============================================================================

\chapter[Prologue : Quarante-Cinq Questions]{Prologue : Quarante-Cinq Questions et Aucune Réponse Fiable}
\label{chap:prologue}
\chaptermark{Prologue}

\epigraph{%
    « Une grille qui interdit de répondre \emph{je ne sais pas}\\
    ne supprime pas l'ignorance~: elle la déguise en certitude. »%
}{Thierry Minka, auditeur des systèmes d'information}

\niveauChapitre{Introductif}{%
    Aucun prérequis. Ce prologue se lit avant toute définition~: il part
    d'une mission qui a échoué.%
}

\begin{boxobjectifs}
\begin{itemize}
    \item Reconstituer la chronologie d'une mission d'audit qui a échoué
          (opération LIZBIZ).
    \item Distinguer une impression, une constatation et une opinion
          d'audit.
    \item Expliquer pourquoi une grille binaire OUI/NON produit des réponses
          incohérentes sur un système partiellement observable.
    \item Identifier les trois états épistémiques d'une réponse~: savoir que
          oui, savoir que non, ne pas savoir.
    \item Énoncer les cinq réflexes de l'auditeur, dont le droit de répondre
          \NEK{} (connaissance insuffisante).
    \item Situer le fil rouge du cours~: la méthode auditologique de conduite
          de la mission.
\end{itemize}
\end{boxobjectifs}

\bigskip

% ============================================================================
\section*{Avant de commencer : un rapport que personne n'a voulu signer}
% ============================================================================

Douala, mardi 14~avril~2026. Le comité de direction de LiZbiZ~SA\footnote{La
société LiZbiZ, ses dirigeants, ses prestataires, les dates et les montants
de ce cas sont fictifs. Les mécanismes d'audit, les normes et les textes
cités sont exacts. Le cas reprend l'illustration de l'article
\emph{Vers l'auditologie}~\cite{minka2026auditologie}, dont les données sont
conservées.} examine le rapport d'audit interne commandé deux mois plus tôt.
LiZbiZ est une entreprise de logistique~: quatre centres de données, à
Douala, Yaoundé, Garoua et Bafoussam, et douze agences régionales qui
suivent les expéditions, gèrent les entrepôts et facturent les clients sur
une même application intégrée.

La mission était née de deux alertes. L'inventaire annuel des entrepôts a
fait apparaître un écart de \textbf{186~millions de francs CFA} entre le
stock enregistré et le stock compté. Le commissaire aux comptes, le
cabinet NDONGO \& ASSOCIÉS, a assorti sa certification d'une réserve sur
les stocks. Le directeur général, M.~ETEME, a demandé un audit «~de tout le
système d'information, et vite~».

L'équipe d'audit a travaillé avec une grille de \textbf{45~questions}
réparties en \textbf{5~domaines}, à cocher OUI ou NON. Les réponses ont été
recueillies auprès du responsable informatique, des chefs d'agence et des
responsables des centres de données. Le rapport conclut que le contrôle
d'accès est «~globalement maîtrisé~».

Le commissaire aux comptes en demande les papiers de travail. Il n'y en a
pas. Le responsable informatique conteste deux conclusions. Trois chefs
d'agence affirment n'avoir jamais été interrogés sur les points qu'on leur
attribue. Le comité de direction refuse de valider le rapport.

\medskip

Ce prologue reconstitue cet échec et en tire les réflexes qui structurent
tout le cours. Il ne suppose aucune connaissance préalable.

\begin{boxalert}
\textbf{Violation typique.} Conclure «~globalement maîtrisé~» sur un domaine
où les répondants se contredisent, sans avoir noté qui a répondu quoi, ni
sur quelle preuve.

\smallskip
\textbf{Règle~:} une conclusion d'audit qui ne peut pas être reliée à une
preuve conservée n'est pas une conclusion~: c'est une opinion personnelle.
\end{boxalert}

\separateur

% ============================================================================
\section{Ce qui s'est passé}
\label{sec:pro-chronologie}
% ============================================================================

\subsection{Reconstituer avant de juger}

Avant de juger une mission, on en reconstitue le déroulement à partir des
traces disponibles~: courriels, versions du rapport, comptes rendus. Le
Tableau~P.1 donne la chronologie établie après coup par la
contrôleuse interne, Mme~NKENG.

\begin{tableaucours}{L{2.6cm}L{9.6cm}}{%
    \tableauHeader{Date} &
    \tableauHeader{Événement}
}
12 février 2026 & Le directeur général demande un audit « de tout le système
d'information » par courriel, sans préciser ni objectifs, ni périmètre, ni
référentiel. \\
16 février & L'équipe d'audit reprend une grille de 45 questions utilisée
dans une autre entreprise, en 5 domaines (A à E). \\
17 février -- 13 mars & Réponses recueillies par téléphone et par courriel.
Personne ne note l'identité des répondants question par question. \\
16 mars & Le domaine D (contrôle d'accès) présente des réponses
contradictoires. L'équipe retient la réponse majoritaire. \\
30 mars & Rapport remis~: contrôle d'accès « globalement maîtrisé ». \\
7 avril & Le commissaire aux comptes demande les papiers de travail~: il n'en
existe pas. \\
14 avril & Le comité de direction refuse de valider le rapport. \\
\end{tableaucours}
\vspace{-0.8em}
\begin{center}{\small\itshape Tableau~P.1. Chronologie de la mission
d'audit de LiZbiZ (reconstitution)}\end{center}

\subsection{Impression, constatation, opinion}

Le rapport mélange trois choses qu'il faut toujours séparer.

\begin{boxdef}
\termecle{Impression}\index{impression}~: jugement formé sans démarche ni
preuve, souvent à partir d'un entretien ou d'une rumeur. Elle oriente
l'investigation~; elle n'entre jamais dans un rapport.

\medskip
\termecle{Constatation}\index{constatation} (\textit{finding})~: fait
vérifiable, établi par l'auditeur et appuyé sur une preuve conservée. Un
tiers doit pouvoir la refaire et aboutir au même résultat.

\medskip
\termecle{Opinion d'audit}\index{opinion d'audit}~: conclusion motivée que
l'auditeur tire d'un ensemble de constatations, au regard d'un critère
connu à l'avance, et dont il assume la responsabilité.
\end{boxdef}

« Le contrôle d'accès est globalement maîtrisé » se présente comme une
opinion. Mais elle ne repose sur aucune constatation conservée, et aucun
critère n'avait été fixé. Ce n'est donc, au mieux, qu'une impression.

\subsection{Trois défauts de méthode}

\begin{enumerate}
    \item \textbf{Pas de mandat.} Un courriel n'est pas une lettre de
          mission~: ni objectifs, ni périmètre, ni référentiel. Personne ne
          sait ce que l'audit devait démontrer (Chapitre~\ref{chap:planifier}).
    \item \textbf{Pas de critère.} La grille vient d'une autre entreprise.
          Ses questions ne disent pas ce qui serait « conforme » chez
          LiZbiZ (Chapitre~\ref{chap:decomposition}).
    \item \textbf{Pas de preuve.} Les réponses ne sont pas rattachées à des
          répondants, ni à des documents, ni à des traces
          (Chapitre~\ref{chap:preuve}).
\end{enumerate}

Ces trois défauts suffiraient à disqualifier le rapport. Mais il y a un
quatrième problème, plus profond~: même conduite avec soin, la grille aurait
échoué sur le domaine~D.

% ============================================================================
\section{Pourquoi la grille a échoué}
\label{sec:pro-grille}
% ============================================================================

\subsection{Des réponses qui se contredisent}

Le Tableau~P.2 reproduit trois questions du domaine~D,
telles que trois interlocuteurs y ont répondu.

\begin{tableaucours}{L{5.6cm}ccc}{%
    \tableauHeader{Question (domaine D)} &
    \tableauHeader{Resp. informatique} &
    \tableauHeader{Chef d'agence} &
    \tableauHeader{Centre de Garoua}
}
Les comptes des agents partis sont-ils désactivés~? & OUI & NON & OUI \\
Les droits d'accès sont-ils revus chaque trimestre~? & OUI & NON & NON \\
Les mots de passe des comptes d'administration sont-ils changés~? & OUI & OUI & NON \\
\end{tableaucours}
\vspace{-0.8em}
\begin{center}{\small\itshape Tableau~P.2. Trois questions du domaine D et
leurs réponses binaires (extrait)}\end{center}

Qui ment~? Probablement personne. Le responsable informatique répond pour la
procédure qu'il a rédigée. Le chef d'agence répond pour ce qu'il voit dans
son agence. Le responsable du centre de Garoua répond pour son site. Et
plusieurs d'entre eux, interrogés plus tard, reconnaissent qu'ils
\emph{ne savaient pas} vraiment~: la grille ne leur laissait que OUI ou NON,
ils ont choisi la réponse qui leur paraissait la plus probable, ou la plus
confortable.

\subsection{Trois états, deux cases}

Une réponse d'audit peut traduire trois états de connaissance~: le
répondant sait que le contrôle est en place, il sait qu'il ne l'est pas, ou
il ne le sait pas. Une grille binaire n'offre que deux cases. Le troisième
état est alors forcé dans l'une des deux, au hasard ou par calcul. C'est
l'observation de départ de l'article qui fonde ce cours~\cite{minka2026nek}.

\begin{boxdef}
\termecle{Réponse NEK}\index{NEK (réponse)} (\textit{Not Enough Knowledge},
connaissance insuffisante)~: troisième modalité de réponse, distincte de
\OUI{} et de \NON{}, qui déclare que l'information disponible ne permet pas
de conclure. Elle n'est pas un aveu d'échec~: c'est une information sur
l'état de connaissance de l'organisation.

\medskip
\termecle{Instrument trichotomique}\index{instrument trichotomique}~: grille
dont chaque question admet les trois réponses \OUI{}, \NON{} et \NEK{}.
\end{boxdef}

Les réponses contradictoires du domaine~D ne sont donc pas du bruit à
éliminer en retenant la majorité. Elles signalent que l'organisation
elle-même ne sait pas si ses contrôles d'accès fonctionnent. C'est
précisément ce qu'un auditeur doit découvrir, et ce que la grille binaire
cachait.

\begin{boiteRemarque}{Retenir la réponse majoritaire, est-ce raisonnable~?}
Non. Une majorité de OUI peut venir de répondants qui décrivent la
procédure écrite sans l'avoir vu appliquer. Une réponse minoritaire peut
venir du seul répondant qui a regardé la réalité. Un audit ne vote pas~: il
cherche la preuve. Le vote masque de plus la question la plus utile~:
\emph{pourquoi} les réponses divergent-elles~?
\end{boiteRemarque}

% ============================================================================
\section{La mission reprise}
\label{sec:pro-reprise}
% ============================================================================

En mai~2026, le comité de direction confie une seconde mission au cabinet
AUDIX CONSEIL. Le chef de mission, M.~BEYALA, conserve le domaine~D comme
cible prioritaire et applique la démarche que ce cours enseigne. Les
résultats décrits ci-dessous sont ceux de l'illustration de l'article
\emph{Vers l'auditologie}~\cite{minka2026auditologie}.

\subsection{Cinq décisions de méthode}

\begin{enumerate}
    \item \textbf{Ajouter la modalité \NEK{}.} Chaque question admet
          désormais trois réponses.
    \item \textbf{Mesurer l'incertitude.} Pour le domaine~D, l'équipe
          calcule l'entropie des réponses~: environ \textbf{1,2~bit}, sur un
          maximum de $\log_2 3 \approx 1{,}585$~bit. C'est le signe d'une
          forte opacité~: l'organisation ne sait pas ce qu'elle fait en
          matière d'accès. Reste à savoir si cette ignorance est sincère ou
          calculée.
    \item \textbf{Décomposer le domaine.} Le domaine~D est découpé en
          $5 \times 5$ sous-éléments, de sorte qu'aucun niveau ne demande à
          l'auditeur de tenir plus de sept éléments à l'esprit à la fois.
    \item \textbf{Comparer les contextes de réponse.} Les réponses données en
          entretien formel, où l'enjeu est élevé, diffèrent nettement de
          celles de l'auto-évaluation, où l'enjeu est faible~: certains
          répondants utilisent \NEK{} de façon stratégique.
    \item \textbf{Regarder dans le temps.} À partir de trois cycles d'audit
          antérieurs, une estimation de l'évolution du domaine~D donne
          \textbf{35~\%} de probabilité de régression dans les 18~mois.
\end{enumerate}

\subsection{Lire un domaine par trois mesures}

Le chapitre~\ref{chap:conclure} définira trois mesures qui accompagnent
tout le cours~: le taux de conformité~$\mu$ (part des \OUI{}), le taux
d'incertitude~$U$ (part des \NEK{}) et l'entropie~$H$ (dispersion des
réponses, en bits). L'encadré suivant les applique au second passage sur le
domaine~D.

\begin{boxlecture}{Domaine D, second passage}
Cinq questions du domaine~D sont reposées à quatre interlocuteurs~: 20
réponses, dont 13~\OUI{}, 2~\NON{} et 5~\NEK{}.

\medskip
\textbf{Constatation.} $\mu = 13/20 = 0{,}65$~; part des \NON{}~$= 0{,}10$~;
$U = 5/20 = 0{,}25$. L'entropie vaut
$H = -(0{,}65\log_2 0{,}65 + 0{,}10\log_2 0{,}10 + 0{,}25\log_2 0{,}25)
\approx 1{,}24$~bit.

\medskip
\textbf{Analyse} (confiance moyenne). Avec 65~\% de \OUI{}, une grille
binaire aurait conclu à un domaine plutôt conforme. Mais un quart des
réponses déclarent ne pas savoir, et l'entropie est élevée~: on ne peut pas
conclure. La conformité apparente est \emph{opaque}~: c'est ce que l'article
\emph{Vers l'auditologie} appelle une \emph{non-conformité inconnue}, qui
appelle d'abord des investigations et une clarification des processus,
avant tout correctif technique (Chapitre~\ref{chap:conclure}).

\medskip
\hfill\lecture{0{,}65}{0{,}25}{1{,}24}
\end{boxlecture}

\begin{boxinfo}
\textbf{Une règle à retenir dès maintenant.} On ne conclut jamais sur une
seule mesure. Une entropie nulle peut signifier une certitude, mais aussi
une ignorance unanime~: si tous les répondants répondent \NEK{}, l'entropie
est nulle et pourtant personne ne sait rien ($U = 1$). Le
chapitre~\ref{chap:conclure} en tirera la règle de décision du cours.
\end{boxinfo}

\subsection{Deux non-conformités, deux remèdes}

Le rapport de la seconde mission distingue deux situations, que la grille
binaire confondait~:

\begin{tableaucours}{L{3.6cm}L{4.0cm}L{4.6cm}}{%
    \tableauHeader{Situation} &
    \tableauHeader{Signature} &
    \tableauHeader{Remède}
}
Non-conformité connue & Entropie faible, non-conformité élevée & Correctifs
techniques, contrôles renforcés \\
Non-conformité inconnue & Entropie élevée & Formation, clarification des
processus, nouvelles investigations \\
\end{tableaucours}
\vspace{-0.8em}
\begin{center}{\small\itshape Tableau~P.3. Deux types de non-conformité
(d'après~\cite{minka2026auditologie})}\end{center}

Six mois plus tard, l'entropie du domaine~D est descendue à environ
\textbf{0,6~bit}~: l'organisation sait désormais ce qu'elle fait en matière
d'accès. Dans le domaine où l'entropie était faible, la conformité est
passée de \textbf{52~\%} à \textbf{84~\%}.

\subsection{Un exemple de décomposition}

Le chapitre~\ref{chap:decomposition} enseignera le gabarit canonique de
décomposition à huit niveaux~: domaine, sous-domaine, processus, activité,
contrôle, test, preuve, trace. Chaque contrôle y porte une \emph{fiche
d'assertion}. En voici une, pour l'une des questions qui posaient problème.

\begin{boxobjet}{D.2.1 · Comptes des agents partis}
\textbf{Position dans le gabarit~:} domaine D (contrôle d'accès)
$\rightarrow$ sous-domaine D.2 (cycle de vie des comptes) $\rightarrow$
processus «~départ d'un agent~» $\rightarrow$ activité «~retrait des
accès~» $\rightarrow$ contrôle D.2.1.

\textbf{Objectif~:} s'assurer qu'aucun agent parti ne conserve d'accès à
l'application intégrée.

\textbf{Critère~:} le compte d'un agent parti est désactivé dans les
48~heures qui suivent la date de départ enregistrée par les ressources
humaines.

\textbf{Question~:} le compte de chaque agent parti au cours de la période
a-t-il été désactivé dans les 48~heures~? Réponse \OUI{}, \NON{} ou \NEK{}.

\textbf{Test~:} rapprocher la liste des départs de la période et la liste
des comptes actifs, puis comparer les dates.

\textbf{Preuve et trace~:} extraction des départs signée par les ressources
humaines~; extraction des comptes actifs, avec son empreinte~; journal des
désactivations de l'application.

\textbf{Risque~:} un agent parti utilise son compte pour consulter ou
modifier des données, par exemple pour masquer un écart de stock.

\textbf{Conséquence~:} pertes non détectées, réserve du commissaire aux
comptes reconduite, manquement à l'obligation de sécuriser les données
personnelles des clients.
\end{boxobjet}

On remarque que la question n'a qu'une seule condition, que la preuve
attendue est nommée avant le test, et que la réponse \NEK{} est prévue~:
si les ressources humaines ne peuvent pas fournir la liste des départs,
l'auditeur ne répond ni OUI ni NON, il le constate.

% ============================================================================
\section{Les cinq réflexes de l'auditeur}
\label{sec:pro-reflexes}
% ============================================================================

L'échec de la première mission et la réussite de la seconde se résument en
cinq réflexes. Chacun fait l'objet d'un chapitre du cours.

\begin{tableaucours}{cL{4.2cm}L{6.0cm}c}{%
    \tableauHeader{N°} &
    \tableauHeader{Réflexe} &
    \tableauHeader{Ce qu'il évite} &
    \tableauHeader{Chapitre}
}
1 & Pas d'audit sans mandat écrit & Une mission sans objectif ni périmètre,
dont on conteste ensuite les conclusions & \ref{chap:planifier} \\
2 & Le critère avant le constat & Juger « maîtrisé » sans savoir ce qui
serait conforme & \ref{chap:decomposition} \\
3 & La preuve avant la conclusion & Des affirmations qu'aucun tiers ne peut
vérifier & \ref{chap:preuve} \\
4 & Le droit de répondre \NEK{} & Des certitudes forcées, des contradictions
tranchées par un vote & \ref{chap:preuve} \\
5 & Le contradictoire avant le rapport & Un audité qui découvre les
conclusions et les conteste & \ref{chap:mission} \\
\end{tableaucours}
\vspace{-0.8em}
\begin{center}{\small\itshape Tableau~P.4. Les cinq réflexes de
l'auditeur}\end{center}

\begin{boxprocedure}
\textbf{Le minimum du papier de travail.} Pour qu'une réponse d'audit puisse
fonder une constatation, le papier de travail qui la porte indique au
moins~:
\begin{itemize}
    \item la question posée et le contrôle auquel elle se rattache~;
    \item l'identité et la fonction du répondant, la date et le mode de
          recueil (entretien, courriel, observation)~;
    \item la réponse \OUI{}, \NON{} ou \NEK{}, et la preuve qui l'appuie
          (document, extraction, observation)~;
    \item la référence de la preuve conservée, et son empreinte s'il s'agit
          d'un fichier~;
    \item le nom de l'auditeur qui a recueilli la réponse.
\end{itemize}
\end{boxprocedure}

\begin{boxmetier}{Auditeur SI}
\textbf{Reprendre une mission contestée}
\begin{enumerate}
    \item Obtenir un mandat écrit~: objectifs, périmètre (les quatre centres
          de données et les douze agences, ou une partie seulement),
          référentiels, livrables, calendrier.
    \item Ne pas réutiliser les réponses de la première mission~: elles ne
          sont rattachées à aucune preuve.
    \item Décomposer d'abord le domaine contesté, puis rédiger une fiche
          d'assertion par contrôle.
    \item Accepter et documenter les réponses \NEK{}~: elles disent où
          chercher.
    \item Présenter les constatations à l'audité avant le rapport, et
          consigner ses réponses.
\end{enumerate}
\end{boxmetier}

\begin{boxmetier}{Commissaire aux comptes}
\textbf{Ce qu'il attend d'un audit des SI}

Le commissaire aux comptes ne certifie pas le système d'information~: il
certifie les comptes. Mais lorsque les stocks sont gérés par une
application, il doit savoir si les contrôles de cette application sont
fiables. Il demande donc les papiers de travail de l'auditeur des SI~: sans
eux, il ne peut pas s'appuyer sur ses conclusions, et maintient sa réserve.
\end{boxmetier}

% ============================================================================
\section{Le cours et son fil rouge}
\label{sec:pro-cours}
% ============================================================================

\subsection{Une méthode, trois cas}

Ce cours enseigne une méthode de conduite de la mission d'audit des
systèmes d'information, présentée intégralement au fil des chapitres. Elle
est l'instanciation opérationnelle, pour le praticien, du programme de
recherche de l'\termecle{auditologie}\index{auditologie}~: une science
appliquée de la réduction de l'incertitude dans l'audit des
systèmes d'information~\cite{minka2026auditologie}. Son instrument de
réponse \OUI{}~/ \NON{}~/ \NEK{} et ses mesures sont ceux de
l'article~\cite{minka2026nek}. La formalisation de la décomposition, du
suivi dans le temps et des comportements stratégiques est en cours de
publication.

Trois organisations fictives accompagnent le cours~:

\begin{tableaucours}{L{3.0cm}L{5.4cm}L{3.8cm}}{%
    \tableauHeader{Cas} &
    \tableauHeader{Situation} &
    \tableauHeader{Thèmes}
}
LiZbiZ (Douala) & Logistique, quatre centres de données, douze agences~;
écart de stock et réserve du commissaire aux comptes & Méthode, contrôle
interne, applications, sécurité \\
BANQUE SANAGA (Douala) & Banque commerciale après une panne de trois jours
de son système central & Contrôle interne bancaire (COBAC), changements,
continuité \\
ANPP (Yaoundé) & Agence publique de paiement, auditée par la Chambre des
comptes & Audit public, conformité, projets \\
\end{tableaucours}
\vspace{-0.8em}
\begin{center}{\small\itshape Tableau~P.5. Les trois cas fil rouge}\end{center}

Les deux derniers cas viennent du cours MSN-2S3-MSI03 (IT Governance, COBIT
et ITIL)~: les étudiants qui le suivent retrouveront les mêmes
organisations, vues cette fois depuis la position de l'auditeur.

\subsection{Le parcours}

Le cours suit le cycle de la mission. La Partie~I pose les fondements~:
l'audit et l'assurance (Chapitre~\ref{chap:audit-auditologie}), le contrôle
interne (Chapitre~\ref{chap:controle-interne}), le risque et la fraude. La
Partie~II enseigne la méthode, de la préparation à la conclusion. La
Partie~III donne le cadre des normes et du droit. La Partie~IV applique la
méthode aux grands domaines d'audit, jusqu'aux études de cas intégratrices
(Chapitre~\ref{chap:cas-integrateurs}). Les Parties~V et~VI ouvrent sur les
outils, l'intelligence artificielle et la recherche
(Chapitre~\ref{chap:auditologie}).

Les dix laboratoires de NDAP SUKULU font pratiquer chaque étape sur les
pièces de LiZbiZ, de la BANQUE SANAGA et de l'ANPP.

\separateur

\begin{boxresume}
\textbf{Ce qu'il faut retenir de ce prologue~:}
\begin{itemize}
    \item \textbf{Un rapport sans papiers de travail n'est pas un rapport}~:
          une conclusion doit pouvoir être reliée à une preuve conservée.
    \item \textbf{Impression, constatation, opinion}~: seule la constatation
          est un fait vérifiable~; l'opinion se fonde sur des constatations
          et un critère connu à l'avance.
    \item \textbf{Une grille binaire force l'ignorance} dans OUI ou NON~:
          les contradictions qui en résultent sont un signal, pas un bruit.
    \item \textbf{La réponse \NEK{}} déclare une connaissance insuffisante~;
          elle dit où chercher.
    \item \textbf{Trois mesures}, $\mu$, $U$ et $H$, lisent un domaine~; on ne
          conclut jamais sur une seule.
    \item \textbf{Cinq réflexes}~: mandat écrit, critère avant constat,
          preuve avant conclusion, droit au \NEK{}, contradictoire avant le
          rapport.
\end{itemize}
\end{boxresume}

\bigskip

\begin{boxexo}[title={\ding{45}\quad Exercice~P.1. Impression, constatation ou opinion~? \niveaubase}]
Classez chacune des phrases suivantes, extraites du premier rapport, et
justifiez.
\begin{enumerate}[(a)]
    \item « Le responsable informatique nous a paru très compétent. »
    \item « Sur les 14 agents partis en 2025, 5 avaient encore un compte actif
          au 31 janvier 2026 (extraction du 2 février, empreinte jointe). »
    \item « La sécurité des accès est satisfaisante. »
    \item « Au regard de la procédure interne (désactivation sous 48~heures),
          le contrôle de désactivation des comptes n'est pas efficace. »
\end{enumerate}
\end{boxexo}

\begin{boxexo}[title={\ding{45}\quad Exercice~P.2. Réécrire trois questions \niveaumoyen}]
Reprenez les trois questions du Tableau~P.2.
\begin{enumerate}[(a)]
    \item Pour chacune, indiquez le critère implicite et réécrivez-la de
          façon qu'elle ne contienne qu'une seule condition vérifiable.
    \item Indiquez, pour chacune, la preuve qui permettrait de répondre
          \OUI{} ou \NON{}.
    \item Décrivez une situation dans laquelle l'auditeur devrait répondre
          \NEK{}, et ce qu'il fait alors.
\end{enumerate}
\end{boxexo}

\begin{boxexo}[title={\ding{45}\quad Exercice~P.3. La note au directeur général \niveauexpert}]
M.~ETEME vous demande pourquoi il devrait payer une seconde mission.
\begin{enumerate}[(a)]
    \item Rédigez une note d'une page qui explique, sans jargon, pourquoi le
          premier rapport ne peut pas être validé.
    \item Expliquez pourquoi 65~\% de réponses \OUI{} sur le domaine~D ne
          permettent pas de conclure que le domaine est maîtrisé.
    \item Proposez le contenu de la lettre de mission de la seconde mission~:
          objectifs, périmètre, référentiels, livrables.
    \item Indiquez ce que le commissaire aux comptes attendra de cette
          seconde mission pour lever sa réserve.
\end{enumerate}
\end{boxexo}

\bigskip

\begin{boxinfo}
\textbf{Transition.} Le prologue a montré une mission qui échoue et une
mission qui réussit. Avant d'apprendre la méthode, il faut définir ce
qu'est un audit, ce qu'il garantit, et ce qu'il ne garantit pas. Le
Chapitre~\ref{chap:audit-auditologie} pose ces fondements et présente
l'audit comme une réduction de l'incertitude.
\end{boxinfo}
